\section{Conclusions}\label{sec:conclusions}

We set out to measure three small QEC decoder kernels on a commodity FPGA card and report their hardware latency and throughput. What we can report instead, honestly, is a synthesis-level characterisation: the decode logic itself is small and fast by any reasonable estimate, and the specific interface required to read a result back without elevated host privileges costs 35 to 70 times the decoder's own pipeline depth. That is a real, measured, reproducible finding, obtained by resynthesising the same source under the same tool and comparing two interface variants directly, and we think it is a more useful contribution than an unverified real-time claim would have been.

Alongside it, this revision produced real Monte Carlo evidence, verified against an independently written reference decoder, that substantially revises this line of work's central logical-error-rate claim for the Shor code: the measured rate under IID-depolarising noise is roughly an order of magnitude higher than previously reported, a gap confirmed well outside statistical uncertainty at \(10^7\) shots. It also produced the first circuit-level, multi-round result associated with this work, for the repetition code, and a corrected scalability analysis showing precisely where LUT-based decoding stops being viable.

No result in this paper is a hardware measurement. Closing that gap, on hardware this revision has identified and partially characterised but not yet exercised, is the immediate next step, and every number in this paper is structured so that a future hardware run can slot directly into the same tables and figures without requiring a rewrite.
